% This file is the actual content of the blueprint.
% It is consumed by both the web and the print build.
% Do not put \begin{document} here.

% Pull in the LaTeX nodes generated from `@[blueprint]`-tagged Lean
% declarations. The library file is rebuilt by `lake build :blueprint`
% and lives under .lake/build/blueprint/library/.
\IfFileExists{../../.lake/build/blueprint/library/LeanForControl.tex}{%
  \input{../../.lake/build/blueprint/library/LeanForControl}%
}{}

\chapter{Linear systems}

\section{Observability and controllability matrices}

For a finite-dimensional linear system
\[
  \dot{x} = A x + B u, \qquad y = C x,
\]
with state matrix $A$, input matrix $B$, and output matrix $C$, this section
collects the (finite-horizon) controllability and observability matrices and
the bridge theorem connecting the textbook observability predicate to the
kernel of the observability matrix.

\subsection{Observability matrix}

\inputleannode{def:observabilityMatrix}

\inputleannode{def:isObservable}

\inputleannode{lem:observabilityMatrix-apply}

\inputleannode{lem:observabilityMatrix-mulVec-apply}

\subsection{Controllability matrix}

\inputleannode{def:controllabilityMatrix}

\inputleannode{lem:controllabilityMatrix-apply}

\inputleannode{def:isControllable}

\subsection{Milestone: observability iff trivial kernel of \texorpdfstring{$\mathcal{O}$}{O}}

\inputleannode{thm:isObservable-iff-ker-trivial}

\subsection{Rank characterizations}

\inputleannode{thm:isObservable-iff-rank}

\inputleannode{thm:isControllable-iff-rank}

\subsection{Hautus observability}

\inputleannode{def:unobservableSubspace}

\inputleannode{thm:unobservable-eq-bot-iff-observable}

\inputleannode{lem:unobservableSubspace-invariant}

\inputleannode{def:hautusObservabilityMatrix}

\inputleannode{lem:hautus-mulVec-eq-zero-iff}

\inputleannode{thm:not-isObservable-implies-hautus-failure}

\inputleannode{thm:hautus-failure-implies-not-isObservable}

\inputleannode{thm:isObservable-iff-hautus}

\subsection{Hautus controllability}

\inputleannode{def:hautusControllabilityMatrix}

\inputleannode{thm:isControllable-iff-hautus}

\section{Reachable subspace and Kalman decomposition}

\inputleannode{def:reachableSubspace}

\inputleannode{lem:mem-reachableSubspace-iff}

\inputleannode{thm:reachableSubspace-eq-top-iff-controllable}

\inputleannode{lem:reachableSubspace-invariant}

\inputleannode{def:kalman-decomposition}

\inputleannode{thm:kalman-subspaces-exist}

\inputleannode{def:kalman-adapted-equivalence}

\inputleannode{thm:kalman-block-matrix-zero-pattern}

\section{Hurwitz stability}

For a real square matrix $A$, Hurwitz stability is expressed through the complex eigenpairs
of its entrywise complexification. The rate-indexed predicate records a strict spectral
margin, and ordinary Hurwitz stability is the zero-rate case.

\inputleannode{def:isHurwitzWithRate}

\inputleannode{def:isHurwitz}

\inputleannode{thm:isHurwitzWithRate-iff-spectral-shift}

\section{Matrix Lyapunov equations and quadratic certificates}

\inputleannode{def:solvesContinuousLyapunovEquation}

\inputleannode{def:matrixQuadratic}

\inputleannode{thm:hurwitz-lyapunov-equation}

\inputleannode{thm:positive-real-eigenvalue-quadratic-certificate}

\section{Peano--Baker series and the state transition matrix}

For a genuinely time-varying state matrix $A(t)$, this section constructs the state transition
matrix $\Phi(t,t_0)$ as the Peano--Baker series and proves Theorem 5.1: $\Phi$ solves the
matrix ODE $\dot\Phi(t,t_0) = A(t)\Phi(t,t_0)$ with $\Phi(t_0,t_0) = I$, and consequently
$x(t) := \Phi(t,t_0)x_0$ solves $\dot x = A(t)x$, $x(t_0) = x_0$.

\inputleannode{def:peanoBakerTerm}

\inputleannode{def:stateTransitionMatrix}

\inputleannode{lem:continuous-peanoBakerTerm}

\inputleannode{lem:norm-peanoBakerTerm-le}

\inputleannode{lem:norm-peanoBakerTerm-le-of-mem}

\inputleannode{lem:summable-peanoBakerTerm}

\inputleannode{lem:continuousOn-stateTransitionMatrix}

\inputleannode{thm:hasDerivAt-stateTransitionMatrix}

\inputleannode{thm:stateTransitionMatrix-self}

\inputleannode{thm:hasDerivAt-stateTransitionMatrix-mulVec}

\inputleannode{lem:isIntegralSolution-stateTransitionMatrix-mulVec}

\inputleannode{thm:stateTransitionMatrix-mulVec-unique}

\chapter{Ordinary differential equations}

\section{Dini derivatives}

\inputleannode{def:diniDerivRight}

\section{Gronwall--Bellman inequality}

\inputleannode{thm:gronwall-bellman}

\section{Continuous dependence on parameters}

\inputleannode{def:isIntegralSolution}

\inputleannode{thm:continuous-dependence-parameters}

\section{Comparison lemma}

\inputleannode{thm:comparison-lemma}

\chapter{Lyapunov stability}

\section{Trajectories, equilibria, and comparison functions}

\inputleannode{def:isTrajectory}

\inputleannode{def:isEquilibrium}

\inputleannode{def:isClassK}

\inputleannode{def:isClassKInfty}

\section{Stability predicates}

\inputleannode{def:lyapunovStable}

\inputleannode{def:localAsymptoticStable}

\inputleannode{def:globalAsymptoticStable}

The finite-forward predicates below quantify over every solution segment on a
compact interval.  This prevents finite escape from making a stability claim
vacuous merely because a non-equilibrium solution is not globally defined.

\inputleannode{def:isForwardTrajectoryOn}

\inputleannode{def:forwardLyapunovStable}

\inputleannode{def:forwardLocallyExponentiallyStable}

\inputleannode{def:forwardUnstable}

\section{Sublevel sets and positive invariance}

\inputleannode{def:sublevelSet}

\inputleannode{def:isPositivelyInvariant}

\inputleannode{lem:isCompact-sublevel-set}

\section{Lyapunov functions}

\inputleannode{def:isLocalLyapunovFunction}

\inputleannode{def:isStrictLocalLyapunovFunction}

\inputleannode{def:isStrictLyapunovFunction}

\inputleannode{def:isAsymptoticLyapunovFunction}

\section{Milestone theorems for autonomous systems}

\inputleannode{thm:lyapunov-stable}

\inputleannode{thm:lyapunov-local-asymptotic-stable}

\inputleannode{thm:lyapunov-asymptotic-stable}

\inputleannode{thm:lyapunov-global-asymptotic-stable}

\section{Lyapunov's indirect method}

\inputleannode{thm:frechet-centered-remainder-bound}

The primary stable-branch theorem uses the finite-forward trajectory API: it gives
uniform local exponential decay on every finite forward solution segment, without
silently excluding solutions that are not globally defined.

\inputleannode{thm:hurwitz-linearization-forward-locally-exponentially-stable}

The following legacy local-asymptotic-stability result is a compatibility corollary
for globally defined trajectories.

\inputleannode{thm:hurwitz-linearization-local-asymptotic-stable}

\inputleannode{thm:exponential-chetaev-forward-unstable}

\inputleannode{thm:positive-real-eigenvalue-forward-unstable}

\section{LaSalle's invariance principle}

\inputleannode{thm:lasalle-invariance-principle}

\inputleannode{thm:lasalle-local-asymptotic-stable}

\inputleannode{thm:lasalle-global-asymptotic-stable}

\section{Lyapunov class \texorpdfstring{$\mathcal{K}$}{K} bounds}

\inputleannode{thm:lyapunov-class-K-bounds}

\section{Class \texorpdfstring{$\mathcal{KL}$}{KL} functions and the Osgood construction}

\inputleannode{def:isClassKL}

\inputleannode{thm:class-KL-osgood}
